\chapter{TP53 Genetic Testing}

\section*{What Is This Test?}
TP53 genetic testing analyzes the TP53 gene located on chromosome 17p13.1 for germline pathogenic variants. TP53 encodes the p53 protein, a 53-kDa transcription factor and master tumor suppressor that regulates the cellular response to DNA damage, hypoxia, oncogene activation, and other cellular stresses. p53 functions as ``the guardian of the genome'': it induces cell cycle arrest at G1/S (via p21/CDKN1A), promotes DNA repair, and, if the damage is irreparable, triggers apoptosis (via BAX, PUMA, and NOXA). Germline pathogenic variants in TP53 cause Li-Fraumeni syndrome (LFS), one of the most penetrant hereditary cancer predisposition syndromes. Individuals with LFS have a >70\% lifetime risk of developing cancer in males and >90\% in females, with a strikingly young median age at first cancer diagnosis (median age 20--25 years). The LFS tumor spectrum includes: soft tissue sarcomas (rhabdomyosarcoma, leiomyosarcoma), osteosarcoma, premenopausal breast cancer (the most common LFS cancer in women, often bilateral and occurring at young age), brain tumors (gliomas, choroid plexus carcinoma, medulloblastoma), adrenocortical carcinoma (the most characteristic pediatric LFS tumor), and acute leukemias. Survivors of a first LFS-associated cancer have a very high risk of developing a second primary malignancy (radiation-induced, compounded by the underlying TP53 defect). TP53 is also the single most frequently mutated gene across all human cancers: somatic TP53 mutations are found in 50\% of sporadic cancers, including virtually all serous ovarian carcinomas, triple-negative breast cancers, and many therapy-related leukemias. Somatic TP53 mutation is a strong adverse prognostic factor in chronic lymphocytic leukemia (CLL), acute myeloid leukemia (AML), myelodysplastic syndromes (MDS), and multiple myeloma.

\section*{Alternative Names}
\begin{itemize}
  \item TP53 Mutation Analysis
  \item p53 Gene Testing
  \item Li-Fraumeni Syndrome Testing
  \item LFS Genetic Testing
  \item p53 Sequencing
  \item Del(17p) Analysis
  \item TP53 IHC
\end{itemize}
\section*{Principle}
Germline TP53 testing is performed by next-generation sequencing (NGS) of all coding exons (exons 2--11) and flanking intronic splice sites of TP53, with deletion/duplication analysis by MLPA. Somatic TP53 testing is performed on tumor tissue (FFPE) or bone marrow aspirate by NGS as part of comprehensive genomic profiling panels. In CLL and AML, del(17p) is detected by FISH and defines high-risk disease; TP53 mutations are detected concurrently by NGS. Immunohistochemistry for p53 protein is a surrogate marker: missense TP53 mutations typically cause p53 protein accumulation (strong diffuse nuclear overexpression on IHC, ++/+++), while truncating (nonsense, frameshift) mutations and deletions produce complete loss of p53 expression (absent/null pattern). Two patterns indicate a mutated TP53: diffuse strong overexpression (mutant pattern) or complete absence (null pattern); the wild-type pattern is scattered weak heterogeneous staining.

\section*{History}
p53 was discovered in 1979 by several groups (David Lane, Arnold Levine, Lloyd Old) as a 53-kDa cellular protein that co-immunoprecipitated with the SV40 large T antigen. It was initially thought to be an oncogene because of its high expression in cancer cells, but in 1989, Bert Vogelstein and colleagues demonstrated that p53 is a tumor suppressor gene: TP53 mutations are concentrated in colorectal cancers and loss of the wild-type allele occurs through 17p deletion (loss of heterozygosity). The seminal 1990 discovery by Frederick Li and Joseph Fraumeni described the familial cancer syndrome (Li-Fraumeni syndrome) and Malkin et al. identified germline TP53 mutations in LFS families in 1990. In 1991--1994, p53 was shown to be a sequence-specific transcription factor that activates p21, a cyclin-dependent kinase inhibitor, mediating G1 arrest. The 2015 revised Chompret criteria for LFS testing were published, refining the clinical indications for germline TP53 testing, and the 2020 National Comprehensive Cancer Network (NCCN) guidelines include comprehensive LFS screening recommendations including whole-body MRI.

\section*{Why This Test Is Ordered}
Germline TP53 testing for Li-Fraumeni syndrome in individuals meeting the classic LFS criteria (proband with a sarcoma diagnosed before age 45 AND a first-degree relative with any cancer before age 45 AND a first- or second-degree relative with any cancer before age 45 or sarcoma at any age) or the revised Chompret criteria. Somatic TP53 testing in CLL to detect TP53 mutation and/or del(17p) defining high-risk disease requiring targeted therapy (BTK inhibitor, BCL2 inhibitor) rather than chemoimmunotherapy. Somatic TP53 testing in AML and MDS to identify high-risk disease with poor response to conventional chemotherapy. Somatic TP53 testing in solid tumors to identify an adverse prognostic marker.

\section*{Primary vs Secondary Test}
Germline TP53 testing is a primary diagnostic test for Li-Fraumeni syndrome. Somatic TP53 testing is an essential secondary prognostic test in CLL, AML, and MDS.

\section*{How This Test Is Performed}
Blood sample in EDTA (lavender-top) for germline testing. Tumor tissue (FFPE) or bone marrow aspirate for somatic testing.

\section*{What Preparation Is Needed}
Pre-test genetic counseling is recommended given the profound implications of a positive LFS diagnosis. No fasting required.

\section*{Contraindications and Precautions}
\begin{itemize}
  \item The primary contraindication is the absence of informed consent and pre-test genetic counseling. Precautions: Because of the very high cancer risk and radiation sensitivity of TP53 carriers, results carry profound implications, and testing should generally not be performed in minors in the absence of clinical suspicion of a childhood-onset Li-Fraumeni syndrome tumor (adrenocortical carcinoma, choroid plexus carcinoma, rhabdomyosarcoma). A variant of uncertain significance (VUS) is not actionable and should not alter management; reclassification over time is expected. A negative germline result in the absence of a known familial variant does not exclude Li-Fraumeni syndrome when clinical criteria are met, because germline TP53 mutations are detected in only approximately 70--80\% of classic Li-Fraumeni syndrome families. Somatic TP53 mutations in CLL, AML, or MDS reflect the disease, not the germline, and do not indicate Li-Fraumeni syndrome unless concurrent germline testing is positive.
\end{itemize}
\section*{Risks}
Minimal physical risk (bruising, pain at the venipuncture site). The greater risks are psychological (anxiety, depression, distress about cancer risk) and the implications for family members, each of whom has a 50\% probability of carrying the same variant. The Genetic Information Nondiscrimination Act (GINA) of 2008 protects against genetic discrimination in health insurance and employment but does not apply to life, disability, or long-term care insurance.

\section*{What Happens After the Test}
Results typically return in 2--4 weeks and are reported as positive (pathogenic or likely pathogenic germline TP53 variant), negative, or VUS. A positive result establishes Li-Fraumeni syndrome and triggers comprehensive NCCN surveillance, including annual whole-body and brain MRI, breast MRI starting at age 20--25 (with discussion of risk-reducing mastectomy), abdominal and pelvic ultrasound, colonoscopy every 2--5 years starting at age 25, and annual dermatologic examination, with avoidance of radiation whenever possible. Post-test genetic counseling is provided, and cascade testing of at-risk relatives is recommended. For somatic TP53 mutations, results guide therapy selection (for example, BTK or BCL2 inhibitors rather than chemoimmunotherapy in CLL).

\section*{Understanding Results}

\subsection*{Pathogenic/Likely Pathogenic Germline TP53 Variant Detected (Li-Fraumeni Syndrome)}
NCCN LFS surveillance: (1) Comprehensive physical exam including neurologic exam every 6--12 months. (2) Whole-body MRI (brain and total body) annually to detect sarcomas and other solid tumors. (3) Breast MRI with contrast annually starting at age 20--25 (or at diagnosis, whichever is earlier), annual mammogram starting at age 30. Discussion of risk-reducing mastectomy. (4) Brain MRI annually. (5) Abdominal/pelvic ultrasound annually. (6) Colonoscopy every 2--5 years starting at age 25. (7) Annual dermatologic exam. (8) Avoidance of therapeutic and diagnostic radiation whenever possible (due to risk of radiation-induced secondary malignancies in TP53 mutation carriers); MRI and ultrasound should be used preferentially. (9) Cascade testing of family members.

\subsection*{Somatic TP53 Mutation/del(17p) in CLL}
TP53 dysfunction (mutation and/or del(17p)) defines high-risk CLL. Treatment: BTK inhibitor (ibrutinib, acalabrutinib, zanubrutinib) or BCL2 inhibitor (venetoclax) plus anti-CD20 antibody-based regimens are preferred. Chemoimmunotherapy (FCR, BR) is contraindicated due to poor efficacy.

\subsection*{Somatic TP53 Mutation in AML/MDS}
TP53-mutated AML/MDS has a very poor prognosis (median OS <6--12 months) with conventional chemotherapy. Allogeneic stem cell transplantation is the only potentially curative option. Hypomethylating agents (azacitidine, decitabine) or clinical trials are standard.

\subsection*{p53 Immunohistochemistry}
Diffuse strong overexpression OR complete absence = mutant pattern (indicates TP53 mutation). Scattered weak staining of varying intensity = wild-type pattern.

\section*{Normal Results}
Germline TP53: No pathogenic variant detected. Somatic TP53 (tumor/bone marrow): TP53 wild-type (no mutation detected). p53 IHC: Wild-type pattern (weak heterogeneous nuclear staining). Normal p53 function is critical for tumor suppression; loss of p53 function is one of the most fundamental events in human carcinogenesis.

\section*{Follow-up Tests}
For germline TP53 mutation carriers: Annual whole-body MRI (NCCN, 2020), breast MRI, brain MRI, abdominal ultrasound, colonoscopy, dermatologic exam. Cascade genetic testing of first-degree relatives. For somatic TP53 mutation in CLL: FISH for del(17p), immunoglobulin heavy-chain variable (IGHV) mutation analysis. For somatic TP53 in solid tumors: Comprehensive genomic profiling (NGS) to identify co-occurring mutations.

\section*{References}
\begin{enumerate}
  \item Mai PL, Best AF, Peters JA, et al. Risks of first and subsequent cancers among TP53 mutation carriers in the NCI LFS cohort. Cancer. 2016;122(23):3673--3681.
  \item NCCN Genetic/Familial High-Risk Assessment: Breast, Ovarian, and Pancreatic. Version 3.2024.
  \item Stengel A, Kern W, Haferlach T, et al. The impact of TP53 mutations and TP53 deletions on survival varies between AML, ALL, MDS, and CLL. Blood Cancer J. 2017;7(7):e588.
  \item Bouaoun L, Sonber D, Ardin M, et al. TP53 variations in human cancers: new lessons from the IARC TP53 database. Hum Mutat. 2016;37(9):865--876.
  \item Villani A, Shore A, Wasserman JD, et al. Biochemical and imaging surveillance in germline TP53 mutation carriers with Li-Fraumeni syndrome. Lancet Oncol. 2016;17(9):1295--1305.
  \item Malkin D, Li FP, Strong LC, et al. Germ line p53 mutations in a familial syndrome of breast cancer, sarcomas, and other neoplasms. Science. 1990;250(4985):1233--1238.
\end{enumerate}

\clearpage
\section*{Regulatory Status and Authorization Date}
This chapter describes a test category, not a specific commercial product. FDA, CDSCO, EU IVDR, and MHRA approval or registration dates therefore cannot be assigned without the assay manufacturer, product name, instrument, and intended use. For laboratory-developed tests, an individual agency approval date may not exist. See the Preface for the interpretation of regulatory status.

\clearpage
